% =====================================================================
%  Chapter 4 — Principle: Transparency (Say What You Do)
%  Map: research/chapter-map.md (Chapter 4)
%  Shape: idea → origin → cost → what it opens up. Ideas, not how-to.
%  Rule: the author is never quoted in body text; short quotes only in
%  footnotes. Full statements live in interviews/.
% =====================================================================
\chaptersubtitle{Say What You Do}
\chapter{Transparency}
\label{ch:effects}

A calculation takes values and produces a value. Everything else a program
does---printing, reading a sensor, sending a message, reserving memory, calling
foreign code---touches the world. Computer scientists call these
\emph{effects}. Most languages let effects happen anywhere, unannounced.
Silica asks that every one be stated, and stated in a particular place: where
the work is done.

\section{Say it where it happens}

The idea divides a program into two kinds of code. Pure code computes and
touches nothing; it can be read, tested and reasoned about in isolation.
Effectful code does things in the world, and in Silica it lives in a distinct
kind of block that announces exactly which kinds of things it does. The compiler
checks the announcement against the work. It also refuses a block that claims
more than it does, because overstating is as misleading as understating.

The intellectual root is the monad, the device Haskell uses to keep effects
apart from pure computation. Silica keeps the separation and drops most of the
ceremony.

Even the shape of those blocks was chosen with readers in mind. Barney first
tried the familiar \code{do} \ldots\ \code{end} form and dropped it after
feedback from an engineer Barney respects. Those words are used, and overused,
in so many other languages that Silica would forever have to explain how its
meaning differed from all the others: extra effort for every human reader, and
a source of confusion for language models.\footnote{Interview, 2026-11-13.} A
fresh form carries no borrowed meanings.

\section{Answer for your own work}

Here Silica parts company with most effect systems. In them, if one piece of
code prints, everything that uses it inherits the obligation, and the
declaration spreads upward through the whole program. In Silica, each piece of
effectful code answers for its own work and no one else's. A part of the
program that merely uses a printing helper says nothing about printing; the
helper already did.

The question this answers is the one a reviewer actually asks. Not \emph{could
anything, anywhere beneath this, possibly touch the world?}, but \emph{where
does this program touch the world?} Local declarations answer it directly,
because each one sits beside the work it describes.\footnote{The beginner book
calls this ``the honesty rule'': \src{silica/docs/learn-silica.md}, line 185.}

The rule was learned by doing it the other way first. Early Silica attached
effects to every function. What changed was a shift in perspective. Silica is a
functional language, but that does not mean effects belong at the level of the
function. A declaration on a whole function is too coarse a sieve: effects pile
up on it, and they muddy a declaration whose job is to say one thing. In Silica,
as Barney puts it, everything should do one thing and do it well. When the
compiler's own source was converted to the new rule, 238 of its 254
function-level effect declarations turned out to describe nothing the function
itself did. They had been decoration.

The same thinking answers the objection a Haskell programmer will raise: that if
effects do not climb the call chain, a function's type no longer reveals whether
calling it touches the world. That climbing was exactly the idea behind putting
effects on functions in the first place, and it is what Barney found wanting.
When every function's declaration must also carry the effects of everything it
calls, the signal is lost in the noise. It is one of the reasons the marker for
foreign code was moved out of declarations altogether and into the names of
modules, where \code{dangerous_} is pure signal (Chapter~\ref{ch:danger}).

\section{Even memory must be declared}

The most unusual of Silica's effects is memory itself. Reserving or growing
storage is treated as an effect like any other. A pure calculation cannot
quietly consume memory, and code that does must say so and say what kind of
memory it uses. In most languages, memory is an invisible service; here it is
part of what a program admits to doing.

The reason goes to the root of the project. Memory safety was, and remains, one
of the foundational reasons for Silica, and in Barney's view all direct use of
memory is unsafe. It must therefore be called out, and it must be intentional.
That means a little deliberate pain. If reserving memory costs nothing to write,
it happens too often and with too little thought, which is precisely the trouble
with treating memory as an invisible service.

\section{What it costs}

Programs say more. Even a beginner's first program must admit where its memory
comes from. And because declarations do not spread upward, knowing that a piece
of code is used tells you nothing about what it does inside; you look where the
work is.

The pain is real, and it is meant to be. Barney draws the line carefully.
Calling something out, plainly and without obscurity, is useful friction: it
makes people think. Making them do busywork is something else, and it drives
them away.
\section{What it opens up}

\begin{itemize}
  \item \textbf{A map of contact with the world.} Every place a program
        touches files, devices, networks, memory or foreign code is marked in
        the source. Security reviews, safety cases and code audits start from
        that map instead of building one.
  \item \textbf{Pure code that is genuinely free.} Code that declares nothing
        can be tested without a world to test it in, reused without fear, and
        reasoned about without tracing what lies beneath.
  \item \textbf{Memory budgets you can see.} On a device with half a megabyte
        of memory, knowing exactly which code reserves storage is the
        difference between a product and a prototype.
  \item \textbf{Tools that understand programs.} An assistant, a language
        model or an analyser can read effects off the page rather than guess
        them from names. When a model leaves out an effect, or declares
        one it does not need, the compiler's error points to exactly where the
        fix belongs, and the model makes it. In Barney's work this happens
        again and again.
  \item \textbf{Effects that can be composed.} Because effectful blocks are
        small and local, several can sit in one function, and one can sit
        inside another. The monad stands out instead of being hidden.
\end{itemize}
